\chapter{Научные основы коннектомного моделирования и измерения ECAP}

\section{Физический стимул, нейронная реакция и защитное поведение Drosophila}
\ResultPending{Будут разделены физический стимул, нейронная реакция, поведение и субъективная боль.}

\section{Коннектомы Drosophila и анатомические ограничения нейродинамических моделей}
\ResultPending{Планируется сопоставить версии коннектома и ограничения, которые анатомические данные накладывают на модели динамики.}

\section{Формирование ECAP и измерительный канал системы стимуляции спинного мозга}
\ResultPending{Будут описаны возбуждение волокон и регистрация составного ответа.}

\section{Методы моделирования нейродинамики, формирования ECAP и переноса представлений}
\ResultPending{Планируется сопоставить методы моделирования нейродинамики и ECAP, предобучения и переноса представлений.}

\section{Сравнительный анализ существующих подходов и выбор методов сравнения}
\ResultPending{Будут определены ближайшие методы сравнения, включая простую модель кривой роста ECAP.}

\section{Требования к исходным данным и экспериментальной проверке}
\ResultPending{Будут заданы требования к данным и независимой проверке.}

\section*{Выводы по главе}
\ResultPending{Выводы будут сформулированы после реализации задач главы.}
